\section{Discussion}
\label{sec:discussion}

\subsection{Meeting the requirements}
The design meets the cost target (RQ1/NFR1) with headroom: $\approx\num{443}$~EUR
for the full head-plus-rail system against a \num{500}~EUR budget, and it is
built entirely from open hardware and 3D-printed parts, reproducible via Ansible
(NFR2). On-device inference on the Pi~5 + Hailo-8 keeps all compute on the
instrument (NFR3). The remaining requirements---per-chip autofocus, repeatable
translation, unattended incubator operation and the detection accuracy
target---depend on hardware and model work that is not yet complete and are
tracked as milestones M1--M3.

\subsection{Design trade-offs}
Replacing the OpenFlexure XY flexure stage with a horizontal tilt plus an
external linear rail trades planar sample scanning for serial multi-chip
scanning, which matches the perfused-bioreactor use case (one head visiting
several chips) better than a single-sample XY stage. It also simplifies the
printed body to a Z-focus-only part. The cost is a larger mechanical envelope
(rail, carriage, NEMA17) and a new motion-control path that no longer uses the
stock Sangaboard.

\subsection{Open decisions}
Several parameters are deliberately left open and marked in the text:
\begin{itemize}
\item the objective tube-length marking and thread (finite $160$ / RMS vs.\
  $\infty$ ICS / M27), which sets the exact magnification and object-space
  sampling (Section~\ref{sec:optical}, M1);
\item the stage driver architecture---all-Python GPIO versus a Sangaboard-hybrid
  (Section~\ref{sec:electronics}, M2);
\item the biological readout and hence the detector family and CVAT label schema
  (Section~\ref{sec:edgeai}, M1);
\item the chip pitch and rack geometry against the physical chips
  (Section~\ref{sec:mechanical}, M3).
\end{itemize}
The repository documentation is currently inconsistent on the IMX477 sensor size
and the system magnification; these are reconciled in
\texttt{docs/THESIS\_PARAMETER\_GAPS.md} and must be fixed before the optical
budget is finalised.

\subsection{Limitations}
The evaluation is at present a design and cost evaluation; the biological
validation (detection accuracy, unattended runs) is future work. The illumination
per station is not yet characterised (wavelength, diffuser, irradiance), and the
full power budget under imaging, inference and slewing has not been measured.
